% Generated draft: requires mathematical and source review.
\documentclass[11pt]{article}
\usepackage[margin=25mm]{geometry}
\usepackage{fontspec}
\setmainfont{Latin Modern Roman}
\usepackage{amsmath,amssymb,mathtools,unicode-math}
\setmathfont{Latin Modern Math}
\usepackage{enumitem,needspace,xurl,hyperref}
\setlength{\parindent}{0pt}
\setlength{\parskip}{5pt}
\setlength{\emergencystretch}{3em}
\newcommand{\N}{\mathbb{N}}
\newcommand{\Z}{\mathbb{Z}}
\newcommand{\Q}{\mathbb{Q}}
\newcommand{\R}{\mathbb{R}}
\newcommand{\C}{\mathbb{C}}
\newcommand{\F}{\mathbb{F}}
\newcommand{\A}{\mathbb{A}}
\newcommand{\PP}{\mathbb{P}}
\newcommand{\E}{\mathbb{E}}
\newcommand{\eps}{\varepsilon}
\newcommand{\dd}{\,\mathrm{d}}
\begin{document}
\section*{UnsolvedMath ID: 30003282}
\textbf{Source problem code:} catalogue entry
\subsection{Definitions and mathematical statement}
\textbf{Terminology and unresolved scope.} A crystal lattice commonly means a connected graph equipped with a free action of a group \(\Gamma\cong\Z^d\), for some positive integer \(d\), by graph automorphisms, with finite quotient graph. This is a conventional interpretation of the terminology, not a definition specified in the supplied record. The record does not specify a lattice, its dimension, an embedding, or a macroscopic spatial domain.

An interacting particle system is a stochastic process whose configurations assign particle states to lattice sites and whose transition rules describe their interactions. A hydrodynamic limit describes convergence, after spatial and temporal rescaling, of suitable empirical conserved quantities to deterministic macroscopic evolution. For conservative dynamics, the gradient condition usually expresses microscopic conserved currents as discrete gradients of local functions; the precise condition depends on the model. The term \emph{non-gradient} indicates failure of the relevant gradient condition, which the supplied record does not define.

\textbf{Source question.} \emph{Can the hydrodynamic limit be derived for a non-gradient interacting particle system on a crystal lattice?}

\textbf{Qualification.} This is a research question rather than a fully specified mathematical proposition. The record supplies no configuration space, transition generator, interaction assumptions, conserved quantities, initial distributions, scaling, boundary conditions, topology of convergence, or candidate macroscopic equation. It also does not specify whether the requested result concerns one particular system, existence of an example, or every system in an intended class. Consequently, no uniquely quantified hydrodynamic-limit assertion can be extracted without additional source information. The terminology above explains the question and does not add hypotheses to it.
\subsection{Short English statement}
The question asks for a rigorous macroscopic evolution law for a non-gradient particle system on a periodic crystal lattice. The supplied record leaves the system and the limiting assertion unspecified, so a precise theorem statement requires additional hypotheses. Its open status is supplied dated curation, not independently verified here.
\subsection{Sources}
\begin{itemize}
\item Supplied local UnsolvedMath registry record, permanent ID \(30003282\), title \emph{Hydrodynamic Limits for Non-Gradient Crystal-Lattice Systems}, source problem code \texttt{catalogue entry}, in the collection \emph{Oberwolfach Reports Open Problems}. The record attributes its provenance to workshop open-problem discussions documented in Oberwolfach Reports; no individual proposer is supplied.
\item Supplied source citation: \emph{Large Scale Stochastic Dynamics (2017)}, Oberwolfach Reports, DOI \texttt{10.4171/owr/2016/54}, pp.\ 3031--3086. \url{https://doi.org/10.4171/owr/2016/54}. The literature-source list also labels this reference \emph{Oberwolfach Report 54 (2016)}. The record gives the proposed year as 2016; these supplied year labels are retained without reconciliation.
\item Supplied literature source labeled \emph{Guan crystal-lattice hydrodynamics}: \url{https://arxiv.org/abs/2004.09158}. No fuller bibliographic information is supplied.
\item Literature assessment dated exactly \texttt{2026-08-21}: the supplied curation labels the problem open and states that hydrodynamic limits are known for important periodic crystal-lattice models, while the general non-gradient crystal-lattice hydrodynamic limit identified in this record remains unresolved. This is supplied dated curation derived from the OWR material, not an independently verified literature assessment. The record expressly calls for independent verification of time-sensitive status and cited claims.
\end{itemize}
\end{document}
